\section{pH}

\subsection{Aktivita}

Bezrozměrná veličina, popisuje reálné chování roztoku. Na rozdíl od ideálního roztoku, se v reálném roztoku částice navzájem ovlivňují. Aktivita jakékoliv čisté látky v kondenzovaném stavu (kapalina nebo pevná látka) je jednotková. Aktivita plynu závisí na jeho parciálním tlaku, obvykle se označuje jako \textbf{fugacita}.

$\mu_i = \mu_i^0 + \textrm{RT} \ln a_i$\\
$\mu_i$ - chemický potenciál, $\mu_i^0$ - standardní chemický potenciál

Ativitu lze vyjádřit jako součin molární koncentrace a aktivitního koeficientu:\\
a $ = \gamma $c

Aktivitní koeficient je úměrný náboji iontů v roztoku a iontové síle roztoku

\subsubsection{Iontová síla roztoku}
$\log \gamma = -0,509\textrm{z}^2\sqrt{\textrm{I}}$\\
I - iontová síla roztoku - popisuje množství iontů v roztoku\\
$I = \frac{1}{2}\sum\limits_{i=0}^n \textrm{c}_i \textrm{z}^2_i$\\
c$_i$ -- molarita; z$_i$ -- náboj; 0,509 -- konstanta pro vodné roztoky při 25~$^\circ$C

\zadani{Jaké jsou aktivity iontů v roztoku \ce{FeCl3} o koncentraci 0,01 M.}

\ce{FeCl3 -> Fe^{3+} + 3 Cl^-}

Koncentrace iontů jsou 0,01 M pro \ce{Fe^{3+}} a 0,03 M pro \ce{Cl^{-}}.

I = 0,5$\times(0,01\times3^2 + 0,03\times(-1)^2) = 0,06$

\ce{Fe^{3+}}:\\
$\log\gamma = -0,509\times9\times\sqrt{0,06} = -1,116\\
\gamma = 10^{-1,116} = 0,0765\\
a = 0,01\times0,0764 = 0,000764$

\ce{Cl^{-}}:\\
$\log\gamma = -0,509\times1\times\sqrt{0,06} = -0,125\\
\gamma = 10^{-0,125} = 0,750\\
a = 0,03\times0,750 = 0,0225$

\odpoved{Aktivita železitých iontů je 0,000764 a chloridů je 0,0225.}
\newpage

\subsection{Iontový součin vody}
\ce{H2O + H2O <=> H3O+ + OH-}

$\textrm{K = }\frac{[\ce{H3O+}][\ce{OH-}]}{[\ce{H2O}]^2}$

K$_w$ = \ce{[H3O+][OH-]} = 10$^{-14}$

pK$_w$ = pH + pOH = 14

\subsection{Vzorce}

\begin{tabular}{ll}
	Silná kyselina & pH = $-\log$ c \\\hline
	Silná zásada & pH = 14 + $\log$ c \\\hline
	Slabá kyselina & pH = $\frac{1}{2}\textrm{p}K_A-\frac{1}{2}\log$ c \\\hline
	Slabá zásada & pH = $14\ +\ \frac{1}{2}\log\textrm{c} - \frac{1}{2}\textrm{p}K_B$ \\\hline
	Sůl slabé k a silné z & pH = $7\ +\ \frac{1}{2}\log\textrm{c} + \frac{1}{2}\textrm{p}K_A$ \\\hline
	Sůl silné k a slabé z & pH = $7\ -\ \frac{1}{2}\log\textrm{c} - \frac{1}{2}\textrm{p}K_B$ \\\hline
	Sůl slabé k a slabé z & pH = $7\ +\ \frac{1}{2}\textrm{p}K_A - \frac{1}{2}\textrm{p}K_B$ \\\hline
	Pufr -- kyselina & pH = $\textrm{p}K_A + \log \frac{[A^-]}{[HA]}$ \\\hline
	Pufr -- zásada & pH = 14 - $\textrm{p}K_B - \log \frac{[B^+]}{[BOH]}$ \\
\end{tabular}

\subsection{Silné kyseliny a zásady}
\zadani{Vypočítej pH kyseliny chlorovodíkové o koncentraci 0,3 M.}

\ce{HCl -> H+ + Cl-}
pH = -log c = -log 0,3 = 0,52\\
\hrule
\textit{Vypočítej pH kyseliny sírové o koncentraci 0,3 M.}
\hrule
\ce{H2SO4 -> 2 H+ + SO$_4^{2-}$}\\
pH = -log c = -log (2\times 0,3) = 0,22

\zadani{Vypočítej pH hydroxidu sodného o koncentraci 0,3 M.}

\ce{NaOH -> Na+ + OH-}

pOH = -log c = -log 0,3 = 0,52\\
pH = 14 - pOH = 14 - 0,52 = 13,48
\hrule

\newpage

\zadani{Vypočítej pH kyseliny chlorovodíkové o koncentraci 1$\times$10$^{-8}$ M.}

Chybný výpočet: pH = -log 1.10$^{-8}$ = 8

Je nutné uvažovat iontový součin vody, pro správný výpočet musíme uvažovat tři podmínky:

\begin{enumerate}
	\item \ce{[Cl^-]} = c
	\item \ce{[H+]} = [\ce{OH-}] + [\ce{Cl-}]
	\item K$_w$ = \ce{[H+]}[\ce{OH-}]
\end{enumerate}

Tím získáme kvadratickou rovnici:\\
$x^2 - cx - K_w = 0\\
\textrm{c} = 1,051\times10^{-7}\\
\textrm{pH} = \textbf{6,978}$

\subsection{Slabé kyseliny a zásady}

\begin{tabular}{|l|r@{,}l|l|r@{,}l|}
	\hline
	Kyselina & \multicolumn{2}{c|}{pKa} & Zásada & \multicolumn{2}{c|}{pKb} \\\hline
	\ce{H3BO3} & 9 & 23 & \ce{NH3} & 4 & 75 \\\hline
	\multirow{ 3}{*}{\ce{H3PO4}} & 2 & 16 & \ce{AgOH} & 3 & 96 \\\cline{2-6}
	& 7 & 21 & \multirow{ 2}{*}{\ce{Ca(OH)2}} & 1 & 40 \\\cline{2-3}\cline{5-6}
	& 12 & 32 & & 2 & 43 \\\hline
	k. octová & 4 & 76 & anilin & 9 & 37 \\\hline
	HF & 3 & 17 & & & \\\hline
\end{tabular}

\zadani{Jaké je pH 0,2 M kyseliny octové, p$K_a$ = 4,76?}

\ce{CH3COOH <=> CH3COO^- + H^+}

$K_a = 10^{-\textrm{pK}_a} = 10^{-4.76} = 0,000017$

$K_a = \frac{[\ce{CH3COO^-}][\ce{H^+}]}{[\ce{CH3COOH}]} = \frac{x.x}{0,2-x}$

Dosadíme za $K_a$ a upravíme získaný výraz, čímž dostaneme kvadratickou rovnici:

$x^2 + 0,000017x - 0,0000034 = 0$

Kvadratickou rovnici vyřešíme pomocí diskriminantu:

$x_{1,2} = \frac{-b \pm \sqrt{D}}{2a} = \frac{-b \pm \sqrt{b^2 - 4ac}}{2a} = \frac{-0,000017 \pm \sqrt{0,000017^2 - 4.1.(-0,0000034)}}{2.1}$

Ze dvou vypočítaných kořenů zvolíme ten kladný, koncentrace nemůže být záporná.

$x = 0,001835$

$\textrm{pH} = -\log[H^+] = -\log0,001835 =$ \textbf{2,736}

\textbf{Zjednodušený výpočet}

$K_a = \frac{[\ce{CH3COO^-}][\ce{H^+}]}{[\ce{CH3COOH}]} = \frac{x.x}{0,2}$

Dosadíme za $K_a$ a upravíme získaný výraz, čímž dostaneme kvadratickou rovnici:

$x^2 - 0,0000034 = 0\\
x = \pm\sqrt{0,0000034} = 0,001844$

pH = -log 0,001844 = \textbf{2,734}

\textbf{Vzorec pro výpočet pH:}

pH = $\frac{1}{2}\textrm{p}K_A-\frac{1}{2}\log$ c = $\frac{1}{2}\times4,76 -\frac{1}{2}\log$ 0,2 = 2,73

\zadani{Jaké je pH 0,25 M roztoku amoniaku}

\ce{NH3 + H2O -> NH$^+_4$ + OH-}

pK$_B$ = 4,76\\
K$_B$ = 1,74$\times$10$^{-5}$

K$_B$ = $\frac{[\ce{NH$^+_4$}][\ce{OH-}]}{[\ce{NH3}]}\ =\ \frac{\ce{x^2}}{\ce{c}}\\
\textrm{x = }\sqrt{\textrm{K}_B\times c}\ =\ \sqrt{1,74\times10^{-5}\times 0,25}\ =\ 2,09\times10^{-3}$\\
pOH = -log $2,09\times10^{-3}$ = 2,68\\
pH = 14 - 2,68 = \textbf{11,32}

\textbf{Vzorec:}\\
pH = $14\ +\ \frac{1}{2}\log\textrm{c} - \frac{1}{2}\textrm{p}K_B$ = 14 + 0,5 log 0,25 - 0,5$\times$4,76 = 11,32

\subsection{Konjugované kyseliny a zásady}
Konjugovaný pár kyselina--zásada vzniká přesunem proton z kyseliny na zásadu.

\ce{H2SO4 + HNO3 <=> HSO$^-_4$ + H2NO$^+_3$\\
CH3COOH + H2O <=> CH3COO^- + H3O+}

\begin{itemize}
	\item Pokud je kyselina silná, je odpovídající konjugovaná zásada slabá.
	\item Pokud je kyselina slabá, je odpovídající konjugovaná zásada silná.
	\item Pokud je kyselina velmi slabá, je odpovídající konjugovaná zásada silná.
\end{itemize}

\subsection{Soli}

\subsubsection{Sůl silné kyseliny a silné zásady}

\ce{NaCl -> Na+ + Cl-}

Při disociaci nedochází ke vzniku \ce{H+}, ani \ce{OH-} iontů, hodnota pH tedy není ovlivněna.

\subsubsection{Sůl silné kyseliny a slabé zásady}
\ce{NH4NO3 + H2O -> NH4^+ + NO3^- + NH3 + H^+}\\
$\textrm{pH} = 7 - \frac{1}{2}(\textrm{pK}_b + \log c)$ \\

\subsubsection{Sůl slabé kyseliny a silné zásady}
\ce{NaF + H2O -> Na^+ + F^- + HF + OH^-}\\
$\textrm{pH} = 7 + \frac{1}{2}(\textrm{pK}_a + \log c)$ \\

\subsubsection{Sůl slabé kyseliny a slabé zásady}
\ce{NH4F + H2O -> NH4^+ + F^- + NH3 + H^+ + HF + OH^-}\\
$\textrm{pH} = 7 + \frac{1}{2}(\textrm{pK}_a - \textrm{pK}_b)$\\
$\textrm{pK}_a (\ce{HF}) = 3,17$\\
$\textrm{pK}_b (\ce{NH3}) = 4,75$\\
$\textrm{pH} = 7 + \frac{1}{2}(3,17 - 4,75) = 6,21$\\

\subsection{Neutralizace}
\zadani{Jaké pH bude mít roztok připravený rozpuštěním 0,853 g KOH v 200 cm$^3$ roztoku \ce{H2SO4} o koncentraci 0,015 M?}

\ce{2 KOH + H2SO4 -> K2SO4 + 2 H2O}

n(KOH) = $\frac{0,853}{56,11}$ = 0,0152 mol\\
n(\ce{H2SO4}) = c.V = 0,015 . 0,20 = 0,0030 mol

Na neutralizaci kyseliny sírové budeme potřebovat dvojnásobné látkové množství KOH, tzn. 0,0060 mol. Po reakci nám zbyde:

n = 0,0152 $-$ 0,0060 = 0,0092 mol KOH\\
c = $\frac{0,0092}{0,2}$ = 0,046 M\\
pH = 14 + log c = 12,66

\odpoved{pH připraveného roztoku bude 12,66.}

\clearpage

\subsection{Pufry}

Jde o směs slabé kyseliny a její soli nebo slabé zásady a její soli. Příkladem je např. acetátový pufr - směs kyseliny octové a octanu sodného.

Rovnováhy v pufru lze popsat rovnicemi:\\
\ce{CH3COOH + H2O <-> CH3COO^- + H3O^+}\\
\ce{CH3COONa + H2O <-> CH3COOH + Na^+ + OH^-}\\
Přídavkem kyseliny vzniknou molekuly kyseliny octové, přídavkem zásady ionty octanu. pH roztoku se nezmění.\\
$\textrm{pH} = \textrm{pK}_a + \log \frac{[A^-]}{[HA]}$\\
$\textrm{pH} = 14 - \textrm{pK}_b + \log \frac{[B]}{[BH^+]}$


\begin{tabular}{|c|c|c|}
	\hline
	\textbf{Pufr} & \textbf{Složení} & \textbf{Rozsah pH} \\\hline
	Acetátový & \ce{CH3COOH}/\ce{CH3COONa} & 3,8 - 5,8 \\\hline
	Fosfátový & \ce{NaH2PO4}/\ce{Na2HPO4} & 6,2 - 8,2 \\\hline
	Borátový & \ce{H3BO3}/\ce{Na2B4O7} & 8,25 - 10,25 \\\hline
\end{tabular}

\zadani{Pufr obsahuje ve 100 cm$^3$ 0,25 molu kyseliny octové a 0,7 molu octanu sodného. Jaké bude jeho pH?}

pK$_a$ = 4,74\\
$[$HA$]$ = $\frac{0,25}{0,1}$ = 2,5 mol.dm$^{-3}$\\
$[$A$^-]$ = $\frac{0,7}{0,1}$ = 7 mol.dm$^{-3}$\\

$\textrm{pH} = \textrm{pK}_a + \log \frac{[A^-]}{[HA]} = 4,74 + \log \frac{7}{2,5}$ = 5,18\\

\odpoved{pH pufru bude 5,18.}

\hrule

\zadani{Pufr obsahuje ve 100 cm$^3$ 0,25 molu amoniaku a 0,7 molu chloridu amonného. Jaké bude jeho pH?}

pK$_b$ = 4,74\\
$[$B$]$ = $\frac{0,25}{0,1}$ = 2,5 mol.dm$^{-3}$\\
$[$BH$^+]$ = $\frac{0,7}{0,1}$ = 7 mol.dm$^{-3}$\\

$\textrm{pH} = 14 - \textrm{pK}_b + \log \frac{[B]}{[BH+]} = 14 - 4,74 + \log \frac{7}{2,5}$ = 9,71\\

\odpoved{pH pufru bude 9,71.}

\clearpage